\subsection{Support through a rank-four automorphic seed}
\label{subsec:rh-rank-four-seed-support}

The ordinary-family theorem constructed the actual ordinary Hecke quotient
in all ranks. Here we use the same quotient, rather than a new ideal
chosen to contain a particular seed. The adjoint vanishing theorem of
Thorne replaces the characteristic-zero adjoint-lifting argument in this
part of the rank-four proof. We also need invariance of the ideal under a
possibly infinite-order anticyclotomic twist.

\paragraph{The precise ordinary datum.}
Keep the Hida setting $s$, its original level $V$ with $V.S'=s.S$, the
coefficient DVR $\mathcal O$, its largeness condition, and the absence of
primitive $\ell$-th roots in $L$. Here $L/F$ is the prescribed CM extension
and $\ell$ is odd. Write $B_0$ for the original universal ordinary
determinant ring, $\Lambda=\mathcal O[[X_1,\ldots,X_N]]$ for its original
weight algebra, and $J_0$ for the kernel of the actual ordinary Hecke
action. We use the established ordinary-family theorem: $B_0/J_0$ is
finite and torsion-free over $\Lambda$; its fixed-weight arithmetic
points with distinct finite parts of the labeled inertial characters
detect $J_0$; the arithmetic occurrence theorem and the finite
anticyclotomic twisting theorem hold for $J_0$. In the rank-four
application $N=4[F:\mathbb Q]$. All assertions below are about that same
$J_0$.

\begin{lemma}[Anticyclotomic twisting of the ordinary ideal]
\label{lem:rh-ordinary-infinite-twist}
Suppose that $t:G_L\to\mathbb Z_\ell$ is continuous, is
anticyclotomic for $L/F$, and vanishes on the inertia groups away from
$\ell$ and on the prescribed local torsion subgroups at $\ell$.
For $u\in 1+\mathfrak m_{\mathcal O}$, put
$\beta(g)=u^{t(g)}$. Twisting by $\beta$ defines an automorphism of
$B_0$ preserving $J_0$.
\end{lemma}
\begin{proof}
The character $(1+Z)^{t(g)}$ over $B_0[[Z]]$ is defined by the binomial
series. It is continuous for the maximal-ideal topology, is trivial on
the prescribed subgroups, is anticyclotomic, and has trivial residue.
Multiply the universal determinant by this character (in degree $i$,
multiply its coefficient at $g$ by $(1+Z)^{it(g)}$) and multiply each
labeled local character by the same character. These are the original
ordinary determinant equations: the polarization is unchanged, each
factor of the ordered ordinary product is scaled by a unit, and the
prescribed inertia and torsion conditions are unchanged. The universal
property therefore gives
$\tau:B_0\to B_0[[Z]]$.

Let $b\in J_0$ and $x$ be an arithmetic point of the actual Hecke
quotient with distinct finite parts of its local characters. The images
of all coefficients of $\tau(b)$ under $x$ are integral and lie in one
finite coefficient field. For every $\ell$-power root of unity $\zeta$,
evaluation at $Z=\zeta-1$ gives the finite-order twist of $x$ by
$\zeta^{t(g)}$. The established finite twisting theorem implies
$x(\tau(b))|_{Z=\zeta-1}=0$. The evaluated power series over the finite
coefficient DVR is zero: a nonzero integral power series has only
finitely many zeros in the open unit disc, by Weierstrass preparation
after removing its maximal common uniformizer power. This argument
does not assume that these roots approach zero.
Fixed-weight density in the ordinary-family theorem now shows,
coefficient by coefficient, that $\tau(b)\in J_0[[Z]]$.
Specializing at $Z=u-1$ converges in $B_0$; ideals in its complete
Noetherian local ring are closed. Consequently the twist sends
$J_0$ into $J_0$. The same proof for $u^{-1}$ gives the inverse.
\end{proof}

If a character is supplied instead by the source's condition
$\ker t\subset\ker\beta$, the image of $t$ is either zero or an open
subgroup $\ell^a\mathbb Z_\ell$ of $\mathbb Z_\ell$. In the nonzero case
divide $t$ by $\ell^a$ and use the image of a generator to obtain the
form of Lemma~\ref{lem:rh-ordinary-infinite-twist}. Its triviality on
the specified subgroups follows from that of $t$. Thus whenever the
original two maps $\phi_0,\phi_n:B_0\to B$ of \texttt{Link.IsPhi}
are defined, with the twists $\kappa$ and $\kappa\beta^n$ respectively,
\[
                       \phi_0(J_0)B=\phi_n(J_0)B .             \tag{RS1}
\]
Indeed the maps agree after precomposition by the twist on the
generators (the determinant coefficients and labeled characters);
the original universal coefficient-generation theorem then gives
their equality as maps. This proves the ideal invariance used by
\texttt{hida\_modular\_ideal\_map}; no construction of $\phi_0$
is being inferred from this identity.

\begin{lemma}[A fixed-weight tangent criterion]
\label{lem:rh-four-tangent-support}
Let $x:B_0\to\overline{\mathbb Q}_\ell$ kill $J_0$, and suppose every
$\mathcal O$-algebra lift of $x$ to dual numbers whose restriction to
$\Lambda$ is constant is itself constant. Then every prime of $B_0$
contained in $\ker x$ contains $J_0$.
\end{lemma}
\begin{proof}
The point is over a finite coefficient field. In fact a local algebra
evaluation from a complete Noetherian local $\mathcal O$-algebra with
residue field finite over that of $\mathcal O$ into
$\overline{\mathbb Q}_\ell$ is automatically integral and continuous:
if a maximal-ideal element had a nonzero image of nonpositive
valuation, a polynomial vanishing on that image could be taken with
integral coefficients and unit constant coefficient (use the minimal
polynomial of its integral inverse). Its evaluation in the original
ring would be a unit. Every element is a residue lift plus a
maximal-ideal element. Finitely many generators of the maximal ideal
and their strictly positive valuations give continuity; approximation
by polynomials in topological generators then puts the whole image in
a finite coefficient field. In particular $B_0/\ker x$ is a finite
$\mathcal O$-module. Any dual-number lift factors through
$B_0/(\ker x)^2$, also finite over $\mathcal O$, since
$\ker x/(\ker x)^2$ is finite over $B_0/\ker x$.
Such lifts are continuous and take values in a finite extension with
dual numbers.

Write $\mathfrak p=\ker x$ and $\mathfrak a=\mathfrak p\cap\Lambda$.
The height of $\mathfrak a$ is $N$ since
$\Lambda/\mathfrak a$ is finite integral over $\mathcal O$ with
characteristic-zero fraction field. The relative tangent hypothesis
makes the relative cotangent space at $\mathfrak p$ zero. More
explicitly the residue fields of the two local rings are finite
separable extensions of $\operatorname{Frac}(\mathcal O)$, so the
dual-number cotangent sequence gives
\[
 \operatorname{edim}(B_{0,\mathfrak p})
          \leq\operatorname{edim}(\Lambda_{\mathfrak a})=N.
\]
The finite torsion-free $\Lambda$-algebra $B_0/J_0$ is nonzero.
Every minimal prime of it contracts to zero (minimal primes are
associated primes and no nonzero element of $\Lambda$ annihilates
a nonzero element). Choose one contained in $\mathfrak p/J_0$.
The corresponding domain is integral over the normal domain
$\Lambda$, and going-down along $\mathfrak a$ gives
$\dim (B_0/J_0)_{\mathfrak p}\geq N$. Thus $B_{0,\mathfrak p}$
is regular local of dimension $N$, hence a domain.
A nonzero ideal in this domain lowers the dimension of its quotient.
It follows that $(J_0)_{\mathfrak p}=0$. This is precisely the claimed
containment in every prime below $\mathfrak p$.
\end{proof}

\begin{theorem}[The ordinary rank-four seed component]
\label{thm:rh-ordinary-four-seed-component}
In the ordinary datum above, suppose that $x$ is represented by a
continuous absolutely irreducible rank-four representation $\rho$ of
$G_L$ with the original symmetric conjugate pairing
\[
        \rho(g)^{\mathsf t}P\rho(cgc^{-1})=\epsilon(g)^{-3}P,
\]
where $\epsilon$ is the $\ell$-adic cyclotomic character and $c$ is
complex conjugation. Suppose $\rho$ is attached to a regular algebraic
conjugate-self-dual cuspidal representation, its restriction to
$G_{L(\mu_{\ell^\infty})}$ is absolutely irreducible, and it is
unramified at the distinguished place. Retain \texttt{Resid.V0ok}:
the residual characters there are unramified, coincide on the
decomposition group, and the residue cyclotomic character is nontrivial.
At every labeled ordinary place suppose the ordered characters on
inertia have the form
\[
         \psi_i=\epsilon^{m_i}\theta_i,\qquad
          m_i-m_{i+1}\geq2,\quad \theta_i\text{ of finite order},
\]
with the $m_i$ allowed to vary by place. The ordered flags are the
ones specified by $x$. Then $x$ kills $J_0$, and every prime of $B_0$
below $\ker x$ contains $J_0$.
\end{theorem}
\begin{proof}
The arithmetic occurrence clause of the ordinary-family theorem
applies: the weights are strictly ordered and the gaps exclude an
equality of two labeled characters differing by $\epsilon$.
Local algebraicity follows from local class field theory and the
displayed inertia formula. The given attachment and pairing are
exactly the conjugate-self-dual automorphic hypothesis. Thus
$J_0\subset\ker x$.

We verify the tangent hypothesis of Lemma~\ref{lem:rh-four-tangent-support}.
Enlarge a finite coefficient field $E$ to contain the point,
representation, and the given dual-number evaluation. We work over
$E[\delta]/(\delta^2)$, where $\delta$ is an infinitesimal parameter.
At the characteristic-zero point, not at the residual point modulo
$\ell$, the determinant is absolutely irreducible. Chenevier's
determinant theorem (Theorem 2.22(i) of \emph{The $p$-adic analytic
space of pseudocharacters of a profinite group}) applied to the
Cayley--Hamilton quotient gives a representation $\rho'$ lifting
$\rho$, unique up to conjugation. It is continuous: choose sixteen
group elements whose matrices form a basis of $M_4(E)$, using
absolute irreducibility and Burnside's theorem. Their lifted matrices
form a basis over the dual numbers, by Nakayama, and coefficients of
each matrix $\rho'(g)$ are obtained from the inverse trace-pairing
matrix applied to its sixteen continuous traces. The same trace
pairing shows that the ordered ordinary product and inertia
relations of the actual determinant ring, not just their
specializations to characteristic zero, are matrix identities for
$\rho'$. In particular $\rho'$ is trivial on every prescribed
residual-trivial old inertia subgroup.

The determinant polarization gives an invertible lifted conjugate
pairing, by the uniqueness statement in the same theorem applied to
the twisted conjugate dual. It can be chosen to reduce to $P$.
Its transpose is also a pairing (apply the relation to $cgc^{-1}$);
averaging with its transpose keeps it invertible and gives a symmetric
pairing. An infinitesimal change of basis makes this pairing equal
to $P$ (over characteristic zero, solve the linearized symmetric
congruence by taking half its discrepancy). Write
$\rho'(g)=(1+\delta z(g))\rho(g)$. It follows that $z$ extends to a
cocycle on $G_F$ with $z(c)=0$ for the polarized adjoint
representation $\operatorname{ad}\rho$: on $G_L$ the action is
conjugation, and on $c$ it is the negative transpose conjugated by
the symmetric pairing. This is the adjoint representation appearing
in Thorne's theorem, not the unpolarized adjoint over $L$.

At an ordinary place choose one inertia element $g_*$ on which the
$\psi_i$ are distinct, using the infinite cyclotomic image on the
common kernel of the finite-order parts. The ordered matrix identity
with all arguments equal to $g_*$ splits the dual-number lattice
into four rank-one eigenspaces, since the four eigenvalue
differences are units over $E[\delta]$. Their order is prescribed.
In that eigenbasis replace only the argument at position $i$ by
$g$. The factors to the left are diagonal and invertible on every
row with index at least $i$; those to the right are diagonal and
invertible on every column with index at most $i$. Thus the entries
in those rows and columns of $\rho'(g)-\psi'_i(g)$ are zero.
For all $i$ this proves the upper triangular flag with diagonal
$\psi'_i$ on the whole decomposition group. The constant-weight
hypothesis makes each $\psi'_i$ equal to $\psi_i$ on inertia.

Pass to a finite extension of the local field killing the finite
inertial factors. The underlying $E$-representation of each
rank-one dual-number graded piece is crystalline: it is the
cyclotomic power $\epsilon^{m_i}$ times an unramified
two-dimensional representation. The same holds for the
one-dimensional graded pieces of $\rho$. We check the extensions,
rather than just their graded pieces. If $U,V$ are crystalline
representations whose cyclotomic exponents in $U$ exceed those in
$V$ by at least two, $W=\operatorname{Hom}(V,U)$ has strictly
positive cyclotomic exponents, all at least two. The local
Bloch--Kato dimension formula gives
$\dim H^1_f(W)=\dim H^0(W)+[K:\mathbb Q_\ell]\dim W$; the local Euler
formula gives the same dimension for $H^1(W)$ because
$H^0(W^*(1))=0$. (The Hodge--Tate weight of $\epsilon$ is $-1$.)
All such extensions are therefore crystalline. Applying this
successively to the ordered filtration proves crystallinity of
both $\rho'$ as an $E$-representation and $\rho$ over this finite
local extension. The resulting self-extension class is in
$H^1_f$ there. Restriction on cohomology with coefficients
$\operatorname{ad}\rho\otimes B_{\rm cris}$ is injective across
a finite local extension in characteristic zero, by
restriction--corestriction. Hence the original class belongs to
$H^1_f$ over the original local field.

At old ramified places away from $\ell$, the inertia condition
already proved for the actual lift implies that the extension
class restricts to zero on a finite-index inertia subgroup.
Again restriction is injective in characteristic zero, so the
class is unramified. At the distinguished place the original
representation is unramified. Write $q$ for the residue cardinality
and $\Phi$ for arithmetic Frobenius. The restriction of $z$ to
inertia is an additive homomorphism; only its pro-$\ell$ tame
quotient can have nonzero image. The relation
$\Phi t\Phi^{-1}=t^q$ gives
$(\operatorname{Ad}\rho(\Phi)-q)z(t)=0$.
All eigenvalues of $\operatorname{Ad}\rho(\Phi)$ reduce to $1$,
since all residual Frobenius eigenvalues coincide, whereas
$q\not\equiv1\bmod\ell$. This operator is invertible and $z(t)=0$.
The other places are unramified by the original quotient.
These checks give the finite local conditions at all places of
$L$, hence at all places of $F$ by restriction to $L$ and
injectivity in characteristic zero.

Thorne's Theorem 6.1 (\emph{On the vanishing of adjoint
Bloch--Kato Selmer groups of irreducible automorphic Galois
representations}, 2022) applies to the cuspidal regular algebraic
conjugate-self-dual attachment and the stated irreducibility on
$G_{L(\mu_{\ell^\infty})}$. It gives
$H^1_f(F,\operatorname{ad}\rho)=0$.
Thus $z$ is a coboundary and $\rho'$ is strictly conjugate to the
constant lift. All determinant coefficients are constant.
The same marked eigenbasis argument at each local $g_*$ now
shows that the labeled characters on the whole decomposition
group are constant as well. The determinant and ordered
characters determine the original coefficient evaluation, by
their topological generation and the continuity already checked.
Every fixed-weight tangent is zero.
Lemma~\ref{lem:rh-four-tangent-support} proves the assertion.
\end{proof}

\paragraph{The seed on the original symplectic space.}
Suppose that the fitted deep tower $T$, its original link $K$,
and the maps $\phi_n$ of \texttt{Link.IsPhi} have been constructed.
Let $y:X.A\to\overline{\mathbb Q}_\ell$ be an \texttt{IsSeedAlt}
point and $R$ its symplectic representation of $G_F$. Write
$(h_i)$ for its squared inertial weights, and $\kappa_s$ for its
seed twist with exponent $a_s$ in the source's squared inertia
formula. The seed condition says that $h_0+a_s$ is even and the
successive gaps of the $h_i$ are even and at least four. At each
ordinary place the link has integral exponent $a$ for $\kappa$
and odd squared exponent $b$ for $\beta$ up to finite characters.
Choose $n=0$ or $1$ with $h_0+n$ even. The point obtained through
$\phi_n$ has representation
\[
                 \rho=R|_{G_L}\,\kappa\beta^n
\]
and ordered inertial exponents $(h_i+2a+nb)/2$, all integral
with gaps at least two. Its ratio twist
$\kappa\beta^n/\kappa_s$ has integral inertia exponents at
every $\ell$-adic place, by the parity condition on $a_s$.
It is algebraic by global class field theory (locally algebraic
at $\ell$, finite on inertia elsewhere). Twisting the seed's
cuspidal regular algebraic attachment by the corresponding
algebraic Hecke character supplies the attachment of $\rho$.
The link identity $\kappa\kappa^c\nu\epsilon^3=1$ and
$\beta\beta^c=1$, with $\nu$ the original similitude, give the
conjugate pairing $P=J R(c)$, where $J$ is the fixed symplectic
form. It is symmetric since $R(c)^2=1$ and $\nu(c)=-1$.
Strong multiplicity one and the paired determinant then give
conjugate self-duality of this cuspidal attachment.
Irreducibility on the cyclotomic kernel is part of
\texttt{IsSeedAlt}. Unramifiedness at the distinguished place is
the original killing condition together with the link's
unramifiedness clause. The link's \texttt{v0} is exactly the
distinguished residual hypothesis of
Theorem~\ref{thm:rh-ordinary-four-seed-component}.

Apply that theorem to the composed evaluation
$x=y|_{X.B}\circ\phi_n$. If $\mathfrak q$ is any prime of
$X.B$ contained in $\ker(y|_{X.B})$, its inverse image under
$\phi_n$ is a prime below $\ker x$, hence contains $J_0$.
By (RS1), $\phi_0(J_0)X.B\subset\mathfrak q$. This proves
the literal \texttt{LocZero} condition for this image ideal.
For two links the conclusion holds for the sum of their image
ideals. This argument establishes the seed clause of the
potential-modularity construction once the original links
and determinant maps are supplied; it does not assume a
rank-four global automorphy propagation theorem. In particular
it uses only the established ordinary \texttt{ModU} for
seed membership, not the stronger \texttt{ModUAlt} local
window conclusion at all arithmetic points.
